\title{2023 年高考数学 新高考II卷}
\date{}
\maketitle

\section*{一、单项选择题}
本大题共8小题，每小题5分，共40分。在每小题给出的四个选项中，只有一项是符合题目要求的。

\begin{question}[points=5]
在复平面内，$(1+3i)(3-i)$对应的点位于（ \quad ）
\begin{choices}
  \item A. 第一象限
  \item B. 第二象限
  \item C. 第三象限
  \item D. 第四象限
\end{choices}
\begin{solution}
\textbf{答案：}A

\textbf{分析：}根据复数的乘法结合复数的几何意义分析判断.

\textbf{详解：}因为$(1+3i)(3-i)=3+8i-3i^2=6+8i$，则所求复数对应的点为$(6,8)$，位于第一象限.
\end{solution}
\end{question}

\begin{question}[points=5]
设集合$A=\{0,-a\}$，$B=\{1,a-2,2a-2\}$，若$A\subseteq B$，则$a=$（ \quad ）
\begin{choices}
  \item A. $2$
  \item B. $1$
  \item C. $\dfrac{2}{3}$
  \item D. $-1$
\end{choices}
\begin{solution}
\textbf{答案：}B

\textbf{分析：}根据包含关系分$a-2=0$和$2a-2=0$两种情况讨论，运算求解即可.

\textbf{详解：}因为$A\subseteq B$，则有：若$a-2=0$，解得$a=2$，此时$A=\{0,-2\}$，$B=\{1,0,2\}$，不符合题意；若$2a-2=0$，解得$a=1$，此时$A=\{0,-1\}$，$B=\{1,-1,0\}$，符合题意；综上所述：$a=1$.
\end{solution}
\end{question}

\begin{question}[points=5]
某学校为了解学生参加体育运动的情况，用比例分配的分层随机抽样方法作抽样调查，拟从初中部和高中部两层共抽取60名学生，已知该校初中部和高中部分别有400名和200名学生，则不同的抽样结果共有（ \quad ）
\begin{choices}
  \item A. $\mathrm{C}_{400}^{45}\cdot\mathrm{C}_{200}^{15}$种
  \item B. $\mathrm{C}_{400}^{20}\cdot\mathrm{C}_{200}^{40}$种
  \item C. $\mathrm{C}_{400}^{30}\cdot\mathrm{C}_{200}^{30}$种
  \item D. $\mathrm{C}_{400}^{40}\cdot\mathrm{C}_{200}^{20}$种
\end{choices}
\begin{solution}
\textbf{答案：}D

\textbf{分析：}利用分层抽样的原理和组合公式即可得到答案.

\textbf{详解：}根据分层抽样的定义知初中部共抽取$60\times\dfrac{400}{600}=40$人，高中部共抽取$60\times\dfrac{200}{600}=20$人，根据组合公式和分步计数原理则不同的抽样结果共有$\mathrm{C}_{400}^{40}\cdot\mathrm{C}_{200}^{20}$种.
\end{solution}
\end{question}

\begin{question}[points=5]
若$f(x)=(x+a)\ln\dfrac{2x-1}{2x+1}$为偶函数，则$a=$（ \quad ）
\begin{choices}
  \item A. $-1$
  \item B. $0$
  \item C. $\dfrac{1}{2}$
  \item D. $1$
\end{choices}
\begin{solution}
\textbf{答案：}B

\textbf{分析：}根据偶函数性质，利用特殊值法求出$a$值，再检验即可.

\textbf{详解：}因为$f(x)$为偶函数，则$f(1)=f(-1)$，即$(1+a)\ln\dfrac{1}{3}=(-1+a)\ln3$，解得$a=0$；当$a=0$时，$f(x)=x\ln\dfrac{2x-1}{2x+1}$，定义域关于原点对称，且$f(-x)=(-x)\ln\dfrac{-2x-1}{-2x+1}=(-x)\ln\dfrac{2x+1}{2x-1}=x\ln\dfrac{2x-1}{2x+1}=f(x)$，故此时$f(x)$为偶函数.
\end{solution}
\end{question}

\begin{question}[points=5]
已知椭圆$C:\dfrac{x^{2}}{3}+y^{2}=1$的左、右焦点分别为$F_{1}$，$F_{2}$，直线$y=x+m$与$C$交于$A$，$B$两点，若$\triangle F_{1}AB$面积是$\triangle F_{2}AB$面积的2倍，则$m=$（ \quad ）
\begin{choices}
  \item A. $\dfrac{2}{3}$
  \item B. $\dfrac{\sqrt{2}}{3}$
  \item C. $-\dfrac{2}{3}$
  \item D. $-\dfrac{\sqrt{2}}{3}$
\end{choices}
\begin{solution}
\textbf{答案：}C

\textbf{分析：}首先联立直线方程与椭圆方程，利用$\Delta>0$，求出$m$范围，再根据三角形面积比得到关于$m$的方程，解出即可.

\textbf{详解：}将直线$y=x+m$与椭圆联立$\begin{cases}y=x+m\\ \dfrac{x^{2}}{3}+y^{2}=1\end{cases}$，消去$y$可得$4x^{2}+6mx+3m^{2}-3=0$，因为直线与椭圆相交于$A,B$点，则$\Delta=36m^{2}-4\times4(3m^{2}-3)>0$，解得$-2<m<2$，设$F_{1}$到$AB$的距离$d_{1}$，$F_{2}$到$AB$的距离$d_{2}$，易知$F_{1}(-\sqrt{2},0),F_{2}(\sqrt{2},0)$，则$d_{1}=\dfrac{|-\sqrt{2}+m|}{\sqrt{2}}$，$d_{2}=\dfrac{|\sqrt{2}+m|}{\sqrt{2}}$，$\dfrac{S_{\triangle F_{1}AB}}{S_{\triangle F_{2}AB}}=\dfrac{|\sqrt{2}+m|}{|\sqrt{2}-m|}=2$，解得$m=-\dfrac{2}{3}$或$-3\sqrt{2}$（舍去）.
\end{solution}
\end{question}

\begin{question}[points=5]
已知函数$f(x)=ae^{x}-\ln x$在区间$(1,2)$上单调递增，则$a$的最小值为（ \quad ）
\begin{choices}
  \item A. $e^{2}$
  \item B. $e$
  \item C. $e^{-1}$
  \item D. $e^{-2}$
\end{choices}
\begin{solution}
\textbf{答案：}C

\textbf{分析：}根据$f'(x)=ae^{x}-\dfrac{1}{x}\ge0$在$(1,2)$上恒成立，再根据分参求最值即可求出.

\textbf{详解：}依题可知，$f'(x)=ae^{x}-\dfrac{1}{x}\ge0$在$(1,2)$上恒成立，显然$a>0$，所以$xe^{x}\ge\dfrac{1}{a}$，设$g(x)=xe^{x},x\in(1,2)$，则$g'(x)=(x+1)e^{x}>0$，所以$g(x)$在$(1,2)$上单调递增，$g(x)>g(1)=e$，故$e\ge\dfrac{1}{a}$，即$a\ge e^{-1}$，即$a$的最小值为$e^{-1}$.
\end{solution}
\end{question}

\begin{question}[points=5]
已知$\alpha$为锐角，$\cos\alpha=\dfrac{1+\sqrt{5}}{4}$，则$\sin\dfrac{\alpha}{2}=$（ \quad ）
\begin{choices}
  \item A. $\dfrac{3-\sqrt{5}}{8}$
  \item B. $\dfrac{-1+\sqrt{5}}{8}$
  \item C. $\dfrac{3-\sqrt{5}}{4}$
  \item D. $\dfrac{-1+\sqrt{5}}{4}$
\end{choices}
\begin{solution}
\textbf{答案：}D

\textbf{分析：}根据二倍角公式（或者半角公式）即可求出.

\textbf{详解：}因为$\cos\alpha=1-2\sin^{2}\dfrac{\alpha}{2}=\dfrac{1+\sqrt{5}}{4}$，而$\alpha$为锐角，解得$\sin\dfrac{\alpha}{2}=\dfrac{3-\sqrt{5}}{8}=\dfrac{(\sqrt{5}-1)^{2}}{16}=\dfrac{\sqrt{5}-1}{4}$.
\end{solution}
\end{question}

\begin{question}[points=5]
记$S_{n}$为等比数列$\{a_{n}\}$的前$n$项和，若$S_{4}=-5$，$S_{6}=21S_{2}$，则$S_{8}=$（ \quad ）
\begin{choices}
  \item A. $120$
  \item B. $85$
  \item C. $-85$
  \item D. $-120$
\end{choices}
\begin{solution}
\textbf{答案：}C

\textbf{分析：}方法一：根据等比数列的前$n$项和公式求出公比，再根据$S_{4},S_{8}$的关系即可解出；方法二：根据等比数列的前$n$项和的性质求解.

\textbf{详解：}方法一：设等比数列$\{a_{n}\}$的公比为$q$，首项为$a_{1}$，若$q=1$则$S_{6}=6a_{1}=3\times2a_{1}=3S_{2}$，与题意不符，所以$q\neq1$；由$S_{4}=-5$，$S_{6}=21S_{2}$可得$\dfrac{a_{1}(1-q^{4})}{1-q}=-5$，$\dfrac{a_{1}(1-q^{6})}{1-q}=21\times\dfrac{a_{1}(1-q^{2})}{1-q}$，由后式可得$1+q^{2}+q^{4}=21$，解得$q^{2}=4$，所以$S_{8}=\dfrac{a_{1}(1-q^{8})}{1-q}=\dfrac{a_{1}(1-q^{4})}{1-q}\times(1+q^{4})=-5\times(1+16)=-85$.
方法二：设等比数列$\{a_{n}\}$的公比为$q$，因为$S_{4}=-5$，$S_{6}=21S_{2}$，所以$q\neq-1$，从而$S_{2},S_{4}-S_{2},S_{6}-S_{4},S_{8}-S_{6}$成等比数列，所以有$(-5-S_{2})^{2}=S_{2}(21S_{2}+5)$，解得$S_{2}=-1$或$S_{2}=\dfrac{5}{4}$；当$S_{2}=-1$时，$S_{2},S_{4}-S_{2},S_{6}-S_{4},S_{8}-S_{6}$即为$-1,-4,-16,S_{8}+21$，易知$S_{8}+21=-64$，即$S_{8}=-85$；当$S_{2}=\dfrac{5}{4}$时，$S_{4}=a_{1}+a_{2}+a_{3}+a_{4}=(a_{1}+a_{2})(1+q^{2})=S_{2}(1+q^{2})>0$，与$S_{4}=-5$矛盾，舍去.
\end{solution}
\end{question}

\section*{二、多项选择题}
本题共4小题，每小题5分，共20分。在每小题给出的选项中，有多项符合题目要求。全部选对的得5分，部分选对的得2分，有选错的得0分。

\begin{question}[points=5]
已知圆锥的顶点为$P$，底面圆心为$O$，$AB$为底面直径，$\angle APB=120^{\circ}$，$PA=2$，点$C$在底面圆周上，且二面角$P-AC-O$为$45^{\circ}$，则（ \quad ）
\begin{choices}
  \item A. 该圆锥的体积为$\pi$
  \item B. 该圆锥的侧面积为$4\sqrt{3}\pi$
  \item C. $AC=2\sqrt{2}$
  \item D. $\triangle PAC$的面积为$\sqrt{3}$
\end{choices}
\begin{solution}
\textbf{答案：}AC

\textbf{分析：}根据圆锥的体积、侧面积判断A、B选项的正确性，利用二面角的知识判断C、D选项的正确性.

\textbf{详解：}依题意，$\angle APB=120^{\circ}$，$PA=2$，所以$OP=1$，$OA=OB=\sqrt{3}$；A选项，圆锥的体积为$\dfrac{1}{3}\times\pi\times(\sqrt{3})^{2}\times1=\pi$，A正确；B选项，圆锥的侧面积为$\pi\times\sqrt{3}\times2=2\sqrt{3}\pi$，B错误；C选项，设$D$是$AC$的中点，连接$OD,PD$，则$AC\perp OD$，$AC\perp PD$，所以$\angle PDO$是二面角$P-AC-O$的平面角，则$\angle PDO=45^{\circ}$，所以$OP=OD=1$，故$AD=CD=\sqrt{3-1}=\sqrt{2}$，则$AC=2\sqrt{2}$，C正确；D选项，$PD=\sqrt{1^{2}+1^{2}}=\sqrt{2}$，所以$S_{\triangle PAC}=\dfrac{1}{2}\times2\sqrt{2}\times\sqrt{2}=2$，D错误.
\end{solution}
\end{question}

\begin{question}[points=5]
设$O$为坐标原点，直线$y=-\sqrt{3}(x-1)$过抛物线$C:y^{2}=2px(p>0)$的焦点，且与$C$交于$M$，$N$两点，$l$为$C$的准线，则（ \quad ）
\begin{choices}
  \item A. $p=2$
  \item B. $|MN|=\dfrac{8}{3}$
  \item C. 以$MN$为直径的圆与$l$相切
  \item D. $\triangle OMN$为等腰三角形
\end{choices}
\begin{solution}
\textbf{答案：}AC

\textbf{分析：}先求得焦点坐标，从而求得$p$，根据弦长公式求得$|MN|$，根据圆与等腰三角形的知识确定正确答案.

\textbf{详解：}A选项：直线$y=-\sqrt{3}(x-1)$过点$(1,0)$，所以抛物线$C:y^{2}=2px(p>0)$的焦点$F(1,0)$，所以$\dfrac{p}{2}=1$，$p=2$，$2p=4$，则A正确，且抛物线$C$的方程为$y^{2}=4x$.
B选项：设$M(x_{1},y_{1})$，$N(x_{2},y_{2})$，由$\begin{cases}y=-\sqrt{3}(x-1)\\ y^{2}=4x\end{cases}$消去$y$并化简得$3x^{2}-10x+3=(x-3)(3x-1)=0$，解得$x_{1}=3$，$x_{2}=\dfrac{1}{3}$，所以$|MN|=x_{1}+x_{2}+p=3+\dfrac{1}{3}+2=\dfrac{16}{3}$，B错误.
C选项：设$MN$的中点为$A$，$M,N,A$到直线$l$的距离分别为$d_{1},d_{2},d$，因为$d=\dfrac{1}{2}(d_{1}+d_{2})=\dfrac{1}{2}(|MF|+|NF|)=\dfrac{1}{2}|MN|$，即$A$到直线$l$的距离等于$|MN|$的一半，所以以$MN$为直径的圆与直线$l$相切，C正确.
D选项：直线$y=-\sqrt{3}(x-1)$即$\sqrt{3}x+y-\sqrt{3}=0$，$O$到直线的距离为$d=\dfrac{\sqrt{3}}{2}$，所以三角形$OMN$的面积为$\dfrac{1}{2}\times\dfrac{16}{3}\times\dfrac{\sqrt{3}}{2}=\dfrac{4\sqrt{3}}{3}$，由上述分析可知$y_{1}=-\sqrt{3}(3-1)=-2\sqrt{3}$，$y_{2}=-\sqrt{3}\left(\dfrac{1}{3}-1\right)=\dfrac{2\sqrt{3}}{3}$，所以$|OM|=\sqrt{3^{2}+(-2\sqrt{3})^{2}}=\sqrt{21}$，$|ON|=\sqrt{\left(\dfrac{1}{3}\right)^{2}+\left(\dfrac{2\sqrt{3}}{3}\right)^{2}}=\dfrac{\sqrt{13}}{3}$，所以三角形$OMN$不是等腰三角形，D错误.
\end{solution}
\end{question}

\begin{question}[points=5]
若函数$f(x)=a\ln x+\dfrac{b}{x}+\dfrac{c}{x^{2}}(a\neq0)$既有极大值也有极小值，则（ \quad ）
\begin{choices}
  \item A. $bc>0$
  \item B. $ab>0$
  \item C. $b^{2}+8ac>0$
  \item D. $ac<0$
\end{choices}
\begin{solution}
\textbf{答案：}BCD

\textbf{分析：}求出函数$f(x)$的导数$f'(x)$，由已知可得$f'(x)$在$(0,+\infty)$上有两个变号零点，转化为一元二次方程有两个不等的正根判断作答.

\textbf{详解：}函数$f(x)=a\ln x+\dfrac{b}{x}+\dfrac{c}{x^{2}}$的定义域为$(0,+\infty)$，求导得$f'(x)=\dfrac{a}{x}-\dfrac{b}{x^{2}}-\dfrac{2c}{x^{3}}=\dfrac{ax^{2}-bx-2c}{x^{3}}$，因为函数$f(x)$既有极大值也有极小值，则函数$f'(x)$在$(0,+\infty)$上有两个变号零点，而$a\neq0$，因此方程$ax^{2}-bx-2c=0$有两个不等的正根$x_{1},x_{2}$，于是$\begin{cases}\Delta=b^{2}+8ac>0\\ x_{1}+x_{2}=\dfrac{b}{a}>0\\ x_{1}x_{2}=-\dfrac{2c}{a}>0\end{cases}$，即有$b^{2}+8ac>0$，$ab>0$，$ac<0$，显然$a^{2}bc<0$，即$bc<0$，A错误，BCD正确.
\end{solution}
\end{question}

\begin{question}[points=5]
在信道内传输$0$，$1$信号，信号的传输相互独立．发送$0$时，收到$1$的概率为$\alpha(0<\alpha<1)$，收到$0$的概率为$1-\alpha$；发送$1$时，收到$0$的概率为$\beta(0<\beta<1)$，收到$1$的概率为$1-\beta$. 考虑两种传输方案：单次传输和三次传输．单次传输是指每个信号只发送$1$次，三次传输是指每个信号重复发送$3$次．收到的信号需要译码，译码规则如下：单次传输时，收到的信号即为译码；三次传输时，收到的信号中出现次数多的即为译码（例如，若依次收到$1$，$0$，$1$，则译码为$1$）.
\begin{choices}
  \item A. 采用单次传输方案，若依次发送$1$，$0$，$1$，则依次收到$1$，$0$，$1$的概率为$(1-\alpha)(1-\beta)^{2}$
  \item B. 采用三次传输方案，若发送$1$，则依次收到$1$，$0$，$1$的概率为$\beta(1-\beta)^{2}$
  \item C. 采用三次传输方案，若发送$1$，则译码为$1$的概率为$\beta(1-\beta)^{2}+(1-\beta)^{3}$
  \item D. 当$0<\alpha<0.5$时，若发送$0$，则采用三次传输方案译码为$0$的概率大于采用单次传输方案译码为$0$的概率
\end{choices}
\begin{solution}
\textbf{答案：}ABD

\textbf{分析：}利用相互独立事件的概率公式计算判断AB；利用相互独立事件及互斥事件的概率计算判断C；求出两种传输方案的概率并作差比较判断D作答.

\textbf{详解：}对于A，依次发送$1$，$0$，$1$，则依次收到$1$，$0$，$1$的事件是发送$1$接收$1$、发送$0$接收$0$、发送$1$接收$1$的3个事件的积，它们相互独立，所以所求概率为$(1-\beta)(1-\alpha)(1-\beta)=(1-\alpha)(1-\beta)^{2}$，A正确.
对于B，三次传输，发送$1$，相当于依次发送$1$，$1$，$1$，则依次收到$1$，$0$，$1$的事件，是发送$1$接收$1$、发送$1$接收$0$、发送$1$接收$1$的3个事件的积，它们相互独立，所以所求概率为$(1-\beta)\cdot\beta\cdot(1-\beta)=\beta(1-\beta)^{2}$，B正确.
对于C，三次传输，发送$1$，则译码为$1$的事件是依次收到$1$，$1$，$0$、$1$，$0$，$1$、$0$，$1$，$1$和$1$，$1$，$1$的事件和，它们互斥，由选项B知，所以所求的概率为$\mathrm{C}_{3}^{2}\beta(1-\beta)^{2}+(1-\beta)^{3}=(1-\beta)^{2}(1+2\beta)$，C错误.
对于D，由选项C知，三次传输，发送$0$，则译码为$0$的概率$P=(1-\alpha)^{2}(1+2\alpha)$，单次传输发送$0$，则译码为$0$的概率$P'=1-\alpha$，而$0<\alpha<0.5$，因此$P-P'=(1-\alpha)^{2}(1+2\alpha)-(1-\alpha)=\alpha(1-\alpha)(1-2\alpha)>0$，即$P>P'$，D正确.
\end{solution}
\end{question}

\section*{三、填空题}
本大题共4小题，每小题5分，共20分。

\begin{question}[points=5]
已知向量$\vec{a}$，$\vec{b}$满足$|\vec{a}-\vec{b}|=\sqrt{3}$，$|\vec{a}+\vec{b}|=|2\vec{a}-\vec{b}|$，则$|\vec{b}|=$ ．
\fillin[{$\sqrt{3}$}]
\begin{solution}
\textbf{答案：}$\sqrt{3}$

\textbf{分析：}法一：根据题意结合向量数量积的运算律运算求解；法二：换元令$\vec{c}=\vec{a}-\vec{b}$，结合数量积的运算律运算求解.

\textbf{详解：}法一：因为$|\vec{a}+\vec{b}|=|2\vec{a}-\vec{b}|$，即$(\vec{a}+\vec{b})^{2}=(2\vec{a}-\vec{b})^{2}$，则$\vec{a}^{2}+2\vec{a}\cdot\vec{b}+\vec{b}^{2}=4\vec{a}^{2}-4\vec{a}\cdot\vec{b}+\vec{b}^{2}$，整理得$\vec{a}^{2}-2\vec{a}\cdot\vec{b}=0$，又因为$|\vec{a}-\vec{b}|=\sqrt{3}$，即$(\vec{a}-\vec{b})^{2}=3$，则$\vec{a}^{2}-2\vec{a}\cdot\vec{b}+\vec{b}^{2}=|\vec{b}|^{2}=3$，所以$|\vec{b}|=\sqrt{3}$.
法二：设$\vec{c}=\vec{a}-\vec{b}$，则$|\vec{c}|=\sqrt{3}$，$\vec{a}+\vec{b}=\vec{c}+2\vec{b}$，$2\vec{a}-\vec{b}=2\vec{c}+\vec{b}$，由题意可得$(\vec{c}+2\vec{b})^{2}=(2\vec{c}+\vec{b})^{2}$，则$\vec{c}^{2}+4\vec{c}\cdot\vec{b}+4\vec{b}^{2}=4\vec{c}^{2}+4\vec{c}\cdot\vec{b}+\vec{b}^{2}$，整理得$\vec{c}^{2}=\vec{b}^{2}$，即$|\vec{b}|=|\vec{c}|=\sqrt{3}$.
\end{solution}
\end{question}

\begin{question}[points=5]
底面边长为4的正四棱锥被平行于其底面的平面所截，截去一个底面边长为2，高为3的正四棱锥，所得棱台的体积为．
\fillin[{$28$}]
\begin{solution}
\textbf{答案：}$28$

\textbf{分析：}方法一：割补法，根据正四棱锥的几何性质以及棱锥体积公式求得正确答案；方法二：根据台体的体积公式直接运算求解.

\textbf{详解：}方法一：由于$\dfrac{2}{4}=\dfrac{1}{2}$，而截去的正四棱锥的高为3，所以原正四棱锥的高为6，所以原正四棱锥的体积为$\dfrac{1}{3}\times(4\times4)\times6=32$，截去的正四棱锥的体积为$\dfrac{1}{3}\times(2\times2)\times3=4$，所以棱台的体积为$32-4=28$.
方法二：棱台的体积为$\dfrac{1}{3}\times3\times(16+4+\sqrt{16\times4})=28$.
\end{solution}
\end{question}

\begin{question}[points=5]
已知直线$l:x-my+1=0$与$\odot C:(x-1)^{2}+y^{2}=4$交于$A$，$B$两点，写出满足“$\triangle ABC$面积为$\dfrac{8}{5}$”的$m$的一个值．
\fillin[{2（$\pm2$或$\pm\dfrac{1}{2}$中任意一个皆可以）}]
\begin{solution}
\textbf{答案：}2（$\pm2$或$\pm\dfrac{1}{2}$中任意一个皆可以）

\textbf{分析：}根据直线与圆的位置关系，求出弦长$AB$，以及点$C$到直线$AB$的距离，结合面积公式即可解出.

\textbf{详解：}设点$C$到直线$AB$的距离为$d$，由弦长公式得$|AB|=2\sqrt{4-d^{2}}$，所以$S_{\triangle ABC}=\dfrac{1}{2}\times d\times2\sqrt{4-d^{2}}=\dfrac{8}{5}$，解得$d=\dfrac{4\sqrt{5}}{5}$或$d=\dfrac{2\sqrt{5}}{5}$，由$d=\dfrac{|1+1|}{\sqrt{1+m^{2}}}=\dfrac{2}{\sqrt{1+m^{2}}}$，所以$\dfrac{2}{\sqrt{1+m^{2}}}=\dfrac{4\sqrt{5}}{5}$或$\dfrac{2}{\sqrt{1+m^{2}}}=\dfrac{2\sqrt{5}}{5}$，解得$m=\pm2$或$m=\pm\dfrac{1}{2}$，故填写其中一个即可.
\end{solution}
\end{question}

\begin{question}[points=5]
已知函数$f(x)=\sin(\omega x+\varphi)$，如图$A$，$B$是直线$y=\dfrac{1}{2}$与曲线$y=f(x)$的两个交点，若$|AB|=\dfrac{\pi}{6}$，则$f(\pi)=$ ．

\begin{center}
  \includegraphics[width=.38\linewidth]{question-16-1.jpg}
\end{center}
\fillin[{$-\dfrac{\sqrt{3}}{2}$}]
\begin{solution}
\textbf{答案：}$-\dfrac{\sqrt{3}}{2}$

\textbf{分析：}设$A\left(x_{1},\dfrac{1}{2}\right),B\left(x_{2},\dfrac{1}{2}\right)$，依题可得$x_{2}-x_{1}=\dfrac{\pi}{6}$，结合$\sin x=\dfrac{1}{2}$的解可得$\omega(x_{2}-x_{1})=\dfrac{2\pi}{3}$，从而得到$\omega$的值，再根据$f\left(\dfrac{2\pi}{3}\right)=0$以及$f(0)<0$，即可得$f(x)=\sin\left(4x-\dfrac{2\pi}{3}\right)$，进而求得$f(\pi)$.

\textbf{详解：}设$A\left(x_{1},\dfrac{1}{2}\right),B\left(x_{2},\dfrac{1}{2}\right)$，由$|AB|=\dfrac{\pi}{6}$可得$x_{2}-x_{1}=\dfrac{\pi}{6}$，由$\sin x=\dfrac{1}{2}$可知，$x=\dfrac{\pi}{6}+2k\pi$或$x=\dfrac{5\pi}{6}+2k\pi$，$k\in\mathbb{Z}$，由图可知，$\omega x_{2}+\varphi-(\omega x_{1}+\varphi)=\dfrac{5\pi}{6}-\dfrac{\pi}{6}=\dfrac{2\pi}{3}$，即$\omega(x_{2}-x_{1})=\dfrac{2\pi}{3}$，$\therefore\omega=4$. 因为$f\left(\dfrac{2\pi}{3}\right)=\sin\left(\dfrac{8\pi}{3}+\varphi\right)=0$，所以$\dfrac{8\pi}{3}+\varphi=k\pi$，即$\varphi=-\dfrac{8\pi}{3}+k\pi$，$k\in\mathbb{Z}$. 所以$f(x)=\sin\left(4x-\dfrac{8\pi}{3}+k\pi\right)=\sin\left(4x-\dfrac{2\pi}{3}+k\pi\right)$，所以$f(x)=\sin\left(4x-\dfrac{2\pi}{3}\right)$或$f(x)=-\sin\left(4x-\dfrac{2\pi}{3}\right)$，又因为$f(0)<0$，所以$f(x)=\sin\left(4x-\dfrac{2\pi}{3}\right)$，$\therefore f(\pi)=\sin\left(4\pi-\dfrac{2\pi}{3}\right)=-\dfrac{\sqrt{3}}{2}$.
\end{solution}
\end{question}

\section*{四、解答题}
本大题共6小题，共70分。解答应写出必要的文字说明、证明过程或演算步骤。

\begin{question}[points=10]
记$\triangle ABC$的内角$A,B,C$的对边分别为$a,b,c$，已知$\triangle ABC$的面积为$\sqrt{3}$，$D$为$BC$中点，且$AD=1$．

（1）若$\angle ADC=\dfrac{\pi}{3}$，求$\tan B$；

（2）若$b^{2}+c^{2}=8$，求$b,c$．
\begin{solution}
\textbf{答案：}（1）$\dfrac{\sqrt{3}}{5}$；（2）$b=c=2$.

\textbf{分析：}（1）方法1，利用三角形面积公式求出$a$，再利用余弦定理求解作答；方法2，利用三角形面积公式求出$a$，作出$BC$边上的高，利用直角三角形求解作答.
（2）方法1，利用余弦定理求出$a$，再利用三角形面积公式求出$\angle ADC$即可求解作答；方法2，利用向量运算律建立关系求出$a$，再利用三角形面积公式求出$\angle ADC$即可求解作答.

\textbf{详解：}（1）方法1：在$\triangle ABC$中，因为$D$为$BC$中点，$\angle ADC=\dfrac{\pi}{3}$，$AD=1$，则$S_{\triangle ADC}=\dfrac{1}{2}AD\cdot DC\sin\angle ADC=\dfrac{1}{2}\times1\times\dfrac{a}{2}\times\dfrac{\sqrt{3}}{2}=\dfrac{\sqrt{3}}{8}a=\dfrac{1}{2}S_{\triangle ABC}=\dfrac{\sqrt{3}}{2}$，解得$a=4$，在$\triangle ABD$中，$\angle ADB=\dfrac{2\pi}{3}$，由余弦定理得$c^{2}=BD^{2}+AD^{2}-2BD\cdot AD\cos\angle ADB$，即$c^{2}=4+1-2\times2\times1\times(-\dfrac{1}{2})=7$，解得$c=\sqrt{7}$，则$\cos B=\dfrac{7+4-1}{2\times\sqrt{7}\times2}=\dfrac{5\sqrt{7}}{14}$，$\sin B=\sqrt{1-\cos^{2}B}=\sqrt{1-\left(\dfrac{5\sqrt{7}}{14}\right)^{2}}=\dfrac{\sqrt{21}}{14}$，所以$\tan B=\dfrac{\sin B}{\cos B}=\dfrac{\sqrt{3}}{5}$.
方法2：在$\triangle ABC$中，因为$D$为$BC$中点，$\angle ADC=\dfrac{\pi}{3}$，$AD=1$，则$S_{\triangle ADC}=\dfrac{1}{2}AD\cdot DC\sin\angle ADC=\dfrac{1}{2}\times1\times\dfrac{a}{2}\times\dfrac{\sqrt{3}}{2}=\dfrac{\sqrt{3}}{8}a=\dfrac{1}{2}S_{\triangle ABC}=\dfrac{\sqrt{3}}{2}$，解得$a=4$，在$\triangle ACD$中，由余弦定理得$b^{2}=CD^{2}+AD^{2}-2CD\cdot AD\cos\angle ADC$，即$b^{2}=4+1-2\times2\times1\times\dfrac{1}{2}=3$，解得$b=\sqrt{3}$，有$AC^{2}+AD^{2}=4=CD^{2}$，则$\angle CAD=\dfrac{\pi}{2}$，$\angle C=\dfrac{\pi}{6}$，过$A$作$AE\perp BC$于$E$，于是$CE=AC\cos C=\dfrac{3}{2}$，$AE=AC\sin C=\dfrac{\sqrt{3}}{2}$，$BE=\dfrac{5}{2}$，所以$\tan B=\dfrac{AE}{BE}=\dfrac{\sqrt{3}}{5}$.
（2）方法1：在$\triangle ABD$与$\triangle ACD$中，由余弦定理得$\begin{cases}c^{2}=\dfrac{1}{4}a^{2}+1-2\times\dfrac{a}{2}\times1\times\cos(\pi-\angle ADC)\\ b^{2}=\dfrac{1}{4}a^{2}+1-2\times\dfrac{a}{2}\times1\times\cos\angle ADC\end{cases}$，整理得$\dfrac{1}{2}a^{2}+2=b^{2}+c^{2}$，而$b^{2}+c^{2}=8$，则$a=2\sqrt{3}$，又$S_{\triangle ADC}=\dfrac{1}{2}\times\sqrt{3}\times1\times\sin\angle ADC=\dfrac{\sqrt{3}}{2}$，解得$\sin\angle ADC=1$，而$0<\angle ADC<\pi$，于是$\angle ADC=\dfrac{\pi}{2}$，所以$b=c=\sqrt{AD^{2}+CD^{2}}=2$.
方法2：在$\triangle ABC$中，因为$D$为$BC$中点，则$2\overrightarrow{AD}=\overrightarrow{AB}+\overrightarrow{AC}$，又$\overrightarrow{CB}=\overrightarrow{AB}-\overrightarrow{AC}$，于是$4|\overrightarrow{AD}|^{2}+|\overrightarrow{CB}|^{2}=(\overrightarrow{AB}+\overrightarrow{AC})^{2}+(\overrightarrow{AB}-\overrightarrow{AC})^{2}=2(b^{2}+c^{2})=16$，即$4+a^{2}=16$，解得$a=2\sqrt{3}$，又$S_{\triangle ADC}=\dfrac{1}{2}\times\sqrt{3}\times1\times\sin\angle ADC=\dfrac{\sqrt{3}}{2}$，解得$\sin\angle ADC=1$，而$0<\angle ADC<\pi$，于是$\angle ADC=\dfrac{\pi}{2}$，所以$b=c=\sqrt{AD^{2}+CD^{2}}=2$.
\end{solution}
\end{question}

\begin{question}[points=10]
$\{a_{n}\}$为等差数列，$b_{n}=\begin{cases}a_{n}-6, & n\text{为奇数}\\ 2a_{n}, & n\text{为偶数}\end{cases}$，记$S_{n}$，$T_{n}$分别为数列$\{a_{n}\}$，$\{b_{n}\}$的前$n$项和，$S_{4}=32$，$T_{3}=16$．

（1）求$\{a_{n}\}$的通项公式；

（2）证明：当$n>5$时，$T_{n}>S_{n}$．
\begin{solution}
\textbf{答案：}（1）$a_{n}=2n+3$；（2）证明见解析.

\textbf{分析：}（1）设等差数列$\{a_{n}\}$的公差为$d$，用$a_{1},d$表示$S_{n}$及$T_{n}$，即可求解作答.
（2）方法1，利用（1）的结论求出$S_{n}$，$b_{n}$，再分奇偶结合分组求和法求出$T_{n}$，并与$S_{n}$作差比较作答；方法2，利用（1）的结论求出$S_{n}$，$b_{n}$，再分奇偶借助等差数列前$n$项和公式求出$T_{n}$，并与$S_{n}$作差比较作答.

\textbf{详解：}（1）设等差数列$\{a_{n}\}$的公差为$d$，而$b_{n}=\begin{cases}a_{n}-6, & n=2k-1\\ 2a_{n}, & n=2k\end{cases}$，$k\in\mathbb{N}^{*}$，则$b_{1}=a_{1}-6$，$b_{2}=2a_{2}=2a_{1}+2d$，$b_{3}=a_{3}-6=a_{1}+2d-6$，于是$\begin{cases}S_{4}=4a_{1}+6d=32\\ T_{3}=4a_{1}+4d-12=16\end{cases}$，解得$a_{1}=5$，$d=2$，$a_{n}=a_{1}+(n-1)d=2n+3$，所以数列$\{a_{n}\}$的通项公式是$a_{n}=2n+3$.
（2）方法1：由（1）知，$S_{n}=\dfrac{n(5+2n+3)}{2}=n^{2}+4n$，$b_{n}=\begin{cases}2n-3, & n=2k-1\\ 4n+6, & n=2k\end{cases}$，$k\in\mathbb{N}^{*}$，当$n$为偶数时，$b_{n-1}+b_{n}=2(n-1)-3+4n+6=6n+1$，$T_{n}=\dfrac{13+(6n+1)}{2}\cdot\dfrac{n}{2}=\dfrac{3}{2}n^{2}+\dfrac{7}{2}n$，当$n>5$时，$T_{n}-S_{n}=\left(\dfrac{3}{2}n^{2}+\dfrac{7}{2}n\right)-(n^{2}+4n)=\dfrac{1}{2}n(n-1)>0$，因此$T_{n}>S_{n}$；当$n$为奇数时，$T_{n}=T_{n+1}-b_{n+1}=\dfrac{3}{2}(n+1)^{2}+\dfrac{7}{2}(n+1)-[4(n+1)+6]=\dfrac{3}{2}n^{2}+\dfrac{5}{2}n-5$，当$n>5$时，$T_{n}-S_{n}=\left(\dfrac{3}{2}n^{2}+\dfrac{5}{2}n-5\right)-(n^{2}+4n)=\dfrac{1}{2}(n+2)(n-5)>0$，因此$T_{n}>S_{n}$，所以当$n>5$时，$T_{n}>S_{n}$.
方法2：由（1）知，$S_{n}=n^{2}+4n$，$b_{n}=\begin{cases}2n-3, & n=2k-1\\ 4n+6, & n=2k\end{cases}$，$k\in\mathbb{N}^{*}$，当$n$为偶数时，$T_{n}=(b_{1}+b_{3}+\cdots+b_{n-1})+(b_{2}+b_{4}+\cdots+b_{n})=\dfrac{-1+2(n-1)-3}{2}\cdot\dfrac{n}{2}+\dfrac{14+4n+6}{2}\cdot\dfrac{n}{2}=\dfrac{3}{2}n^{2}+\dfrac{7}{2}n$，当$n>5$时，$T_{n}-S_{n}=\dfrac{1}{2}n(n-1)>0$，因此$T_{n}>S_{n}$；当$n$为奇数时，若$n\ge3$，则$T_{n}=(b_{1}+b_{3}+\cdots+b_{n})+(b_{2}+b_{4}+\cdots+b_{n-1})=\dfrac{-1+2n-3}{2}\cdot\dfrac{n+1}{2}+\dfrac{14+4(n-1)+6}{2}\cdot\dfrac{n-1}{2}=\dfrac{3}{2}n^{2}+\dfrac{5}{2}n-5$，显然$T_{1}=b_{1}=-1$满足上式，因此当$n$为奇数时，$T_{n}=\dfrac{3}{2}n^{2}+\dfrac{5}{2}n-5$，当$n>5$时，$T_{n}-S_{n}=\dfrac{1}{2}(n+2)(n-5)>0$，因此$T_{n}>S_{n}$，所以当$n>5$时，$T_{n}>S_{n}$.
\end{solution}
\end{question}

\begin{question}[points=10]
某研究小组经过研究发现某种疾病的患病者与未患病者的某项医学指标有明显差异，经过大量调查，得到如下的患病者和未患病者该指标的频率分布直方图：



利用该指标制定一个检测标准，需要确定临界值$c$，将该指标大于$c$的人判定为阳性，小于或等于$c$的人判定为阴性．此检测标准的漏诊率是将患病者判定为阴性的概率，记为$p(c)$；误诊率是将未患病者判定为阳性的概率，记为$q(c)$．假设数据在组内均匀分布，以事件发生的频率作为相应事件发生的概率．

（1）当漏诊率$p(c)=0.5\%$时，求临界值$c$和误诊率$q(c)$；

（2）设函数$f(c)=p(c)+q(c)$，当$c\in[95,105]$时，求$f(c)$的解析式，并求$f(c)$在区间$[95,105]$的最小值．
\begin{solution}
\textbf{答案：}（1）$c=97.5$，$q(c)=3.5\%$；
（2）$f(c)=\begin{cases}-0.008c+0.82, & 95\le c\le100\\ 0.01c-0.98, & 100<c\le105\end{cases}$，最小值为$0.02$.

\textbf{分析：}（1）根据题意由第一个图可先求出$c$，再根据第二个图求出$c\ge97.5$的矩形面积即可解出；
（2）根据题意确定分段点$100$，即可得出$f(c)$的解析式，再根据分段函数的最值求法即可解出.

\textbf{详解：}（1）依题可知，左边图形第一个小矩形的面积为$5\times0.002>0.5\%$，所以$95<c<100$，所以$(c-95)\times0.002=0.5\%$，解得$c=97.5$，$q(c)=0.01\times(97.5-95)+5\times0.002=0.035=3.5\%$.
（2）当$c\in[95,100]$时，$f(c)=p(c)+q(c)=(c-95)\times0.002+(100-c)\times0.01+5\times0.002=-0.008c+0.82\ge0.02$；当$c\in(100,105]$时，$f(c)=p(c)+q(c)=5\times0.002+(c-100)\times0.012+(105-c)\times0.002=0.01c-0.98>0.02$，故$f(c)=\begin{cases}-0.008c+0.82, & 95\le c\le100\\ 0.01c-0.98, & 100<c\le105\end{cases}$，所以$f(c)$在区间$[95,105]$的最小值为$0.02$.
\end{solution}
\end{question}

\begin{question}[points=10]
如图，三棱锥$A-BCD$中，$DA=DB=DC$，$BD\perp CD$，$\angle ADB=\angle ADC=60^{\circ}$，$E$为$BC$的中点．

（1）证明：$BC\perp DA$；

（2）点$F$满足$\overrightarrow{EF}=\overrightarrow{DA}$，求二面角$D-AB-F$的正弦值．

\begin{center}
  \includegraphics[width=.38\linewidth]{question-20-1.jpg}
\end{center}
\begin{solution}
\textbf{答案：}（1）证明见解析；（2）$\dfrac{\sqrt{3}}{3}$.

\textbf{分析：}（1）根据题意易证$BC\perp$平面$ADE$，从而证得$BC\perp DA$；
（2）由题可证$AE\perp$平面$BCD$，所以以点$E$为原点，$ED,EB,EA$所在直线分别为$x,y,z$轴，建立空间直角坐标系，再求出平面$ABD$，$ABF$的一个法向量，根据二面角的向量公式以及同角三角函数关系即可解出.

\textbf{详解：}（1）连接$AE,DE$，因为$E$为$BC$中点，$DB=DC$，所以$DE\perp BC$①，因为$DA=DB=DC$，$\angle ADB=\angle ADC=60^{\circ}$，所以$\triangle ACD$与$\triangle ABD$均为等边三角形，$\therefore AC=AB$，从而$AE\perp BC$②，由①②，$AE\cap DE=E$，$AE,DE\subset$平面$ADE$，所以$BC\perp$平面$ADE$，而$AD\subset$平面$ADE$，所以$BC\perp DA$.
（2）不妨设$DA=DB=DC=2$，$\because BD\perp CD$，$\therefore BC=2\sqrt{2}$，$DE=AE=\sqrt{2}$，$\because AE^{2}+DE^{2}=4=AD^{2}$，$\therefore AE\perp DE$，又$\because AE\perp BC$，$DE\cap BC=E$，$DE,BC\subset$平面$BCD$，$\therefore AE\perp$平面$BCD$．以点$E$为原点，$ED,EB,EA$所在直线分别为$x,y,z$轴，建立空间直角坐标系，如图所示：
设$D(\sqrt{2},0,0)$，$A(0,0,\sqrt{2})$，$B(0,\sqrt{2},0)$，$E(0,0,0)$，设平面$DAB$与平面$ABF$的一个法向量分别为$\vec{n_{1}}=(x_{1},y_{1},z_{1})$，$\vec{n_{2}}=(x_{2},y_{2},z_{2})$，二面角$D-AB-F$平面角为$\theta$，而$\overrightarrow{AB}=(0,\sqrt{2},-\sqrt{2})$，因为$\overrightarrow{EF}=\overrightarrow{DA}=(-\sqrt{2},0,\sqrt{2})$，所以$F(-\sqrt{2},0,\sqrt{2})$，即有$\overrightarrow{AF}=(-\sqrt{2},0,0)$，$\therefore\begin{cases}-\sqrt{2}x_{1}+\sqrt{2}z_{1}=0\\ \sqrt{2}y_{1}-\sqrt{2}z_{1}=0\end{cases}$，取$x_{1}=1$，所以$\vec{n_{1}}=(1,1,1)$；$\begin{cases}\sqrt{2}y_{2}-\sqrt{2}z_{2}=0\\ -\sqrt{2}x_{2}=0\end{cases}$，取$y_{2}=1$，所以$\vec{n_{2}}=(0,1,1)$，所以$\cos\theta=\dfrac{\vec{n_{1}}\cdot\vec{n_{2}}}{|\vec{n_{1}}||\vec{n_{2}}|}=\dfrac{2}{\sqrt{3}\times\sqrt{2}}=\dfrac{\sqrt{6}}{3}$，从而$\sin\theta=\sqrt{1-\dfrac{6}{9}}=\dfrac{\sqrt{3}}{3}$，所以二面角$D-AB-F$的正弦值为$\dfrac{\sqrt{3}}{3}$.
\end{solution}
\end{question}

\begin{question}[points=10]
已知双曲线$C$的中心为坐标原点，左焦点为$(-2\sqrt{5},0)$，离心率为$\sqrt{5}$．

（1）求$C$的方程；

（2）记$C$的左、右顶点分别为$A_{1}$，$A_{2}$，过点$(-4,0)$的直线与$C$的左支交于$M$，$N$两点，$M$在第二象限，直线$MA_{1}$与$NA_{2}$交于点$P$．证明:点$P$在定直线上.
\begin{solution}
\textbf{答案：}（1）$\dfrac{x^{2}}{4}-\dfrac{y^{2}}{16}=1$；（2）证明见解析.

\textbf{分析：}（1）由题意求得$a,b$的值即可确定双曲线方程；
（2）设出直线方程，与双曲线方程联立，然后由点的坐标分别写出直线$MA_{1}$与$NA_{2}$的方程，联立直线方程，消去$y$，结合韦达定理计算可得$\dfrac{x+2}{x-2}=-\dfrac{1}{3}$，即交点的横坐标为定值，据此可证得点$P$在定直线$x=-1$上.

\textbf{详解：}（1）设双曲线方程为$\dfrac{x^{2}}{a^{2}}-\dfrac{y^{2}}{b^{2}}=1(a>0,b>0)$，由焦点坐标可知$c=2\sqrt{5}$，则由$e=\dfrac{c}{a}=\sqrt{5}$可得$a=2$，$b=\sqrt{c^{2}-a^{2}}=4$，双曲线方程为$\dfrac{x^{2}}{4}-\dfrac{y^{2}}{16}=1$.
（2）由(1)可得$A_{1}(-2,0)$，$A_{2}(2,0)$，设$M(x_{1},y_{1})$，$N(x_{2},y_{2})$，显然直线的斜率不为$0$，所以设直线$MN$的方程为$x=my-4$，且$-\dfrac{1}{2}<m<\dfrac{1}{2}$，与$\dfrac{x^{2}}{4}-\dfrac{y^{2}}{16}=1$联立可得$(4m^{2}-1)y^{2}-32my+48=0$，且$\Delta=64(4m^{2}+3)>0$，则$y_{1}+y_{2}=\dfrac{32m}{4m^{2}-1}$，$y_{1}y_{2}=\dfrac{48}{4m^{2}-1}$，直线$MA_{1}$的方程为$y=\dfrac{y_{1}}{x_{1}+2}(x+2)$，直线$NA_{2}$的方程为$y=\dfrac{y_{2}}{x_{2}-2}(x-2)$，联立直线$MA_{1}$与直线$NA_{2}$的方程可得：$\dfrac{x+2}{x-2}=\dfrac{y_{2}(x_{1}+2)}{y_{1}(x_{2}-2)}=\dfrac{y_{2}(my_{1}-2)}{y_{1}(my_{2}-6)}=\dfrac{my_{1}y_{2}-2y_{1}}{my_{1}y_{2}-6y_{1}}=\dfrac{m\cdot\dfrac{48}{4m^{2}-1}-2y_{1}}{m\cdot\dfrac{48}{4m^{2}-1}-6y_{1}}=\dfrac{\dfrac{48m}{4m^{2}-1}-2y_{1}}{\dfrac{48m}{4m^{2}-1}-6y_{1}}=\dfrac{\dfrac{48m}{4m^{2}-1}-2y_{1}}{\dfrac{48m}{4m^{2}-1}-6y_{1}}$，由$y_{1}+y_{2}=\dfrac{32m}{4m^{2}-1}$可得$\dfrac{48m}{4m^{2}-1}=2\cdot\dfrac{16m}{4m^{2}-1}+\dfrac{16m}{4m^{2}-1}=\dfrac{1}{2}(y_{1}+y_{2})+\dfrac{16m}{4m^{2}-1}$，代入得$\dfrac{x+2}{x-2}=\dfrac{\dfrac{1}{2}(y_{1}+y_{2})+\dfrac{16m}{4m^{2}-1}-2y_{1}}{\dfrac{1}{2}(y_{1}+y_{2})+\dfrac{16m}{4m^{2}-1}-6y_{1}}=\dfrac{-\dfrac{3}{2}y_{1}+\dfrac{1}{2}y_{2}+\dfrac{16m}{4m^{2}-1}}{-\dfrac{11}{2}y_{1}+\dfrac{1}{2}y_{2}+\dfrac{16m}{4m^{2}-1}}$，又$y_{2}=\dfrac{32m}{4m^{2}-1}-y_{1}$，代入得$\dfrac{x+2}{x-2}=\dfrac{-\dfrac{3}{2}y_{1}+\dfrac{1}{2}(\dfrac{32m}{4m^{2}-1}-y_{1})+\dfrac{16m}{4m^{2}-1}}{-\dfrac{11}{2}y_{1}+\dfrac{1}{2}(\dfrac{32m}{4m^{2}-1}-y_{1})+\dfrac{16m}{4m^{2}-1}}=\dfrac{-2y_{1}+\dfrac{32m}{4m^{2}-1}}{-6y_{1}+\dfrac{32m}{4m^{2}-1}}=\dfrac{2(-y_{1}+\dfrac{16m}{4m^{2}-1})}{6(-y_{1}+\dfrac{16m}{4m^{2}-1})}=\dfrac{1}{3}$，即$\dfrac{x+2}{x-2}=\dfrac{1}{3}$，解得$x=-4$，故点$P$在定直线$x=-4$上.
\end{solution}
\end{question}

\begin{question}[points=10]
（1）证明：当$0<x<1$时，$x-x^{2}<\sin x<x$；

（2）已知函数$f(x)=\cos ax-\ln(1-x^{2})$，若$x=0$是$f(x)$的极大值点，求$a$的取值范围．
\begin{solution}
\textbf{答案：}（1）证明见详解；（2）$(-\infty,-\sqrt{2})\cup(\sqrt{2},+\infty)$.

\textbf{分析：}（1）分别构建$F(x)=x-\sin x$，$x\in(0,1)$，$G(x)=x-x^{2}+\sin x$，$x\in(0,1)$，求导，利用导数判断原函数的单调性，进而可得结果；
（2）根据题意结合偶函数的性质可知只需要研究$f(x)$在$(0,1)$上的单调性，求导，分类讨论$0<a^{2}\le 2$和$a^{2}\ge 2$，结合（1）中的结论放缩，根据极大值的定义分析求解.

\textbf{详解：}（1）构建$F(x)=x-\sin x$，$x\in(0,1)$，则$F'(x)=1-\cos x>0$对$\forall x\in(0,1)$恒成立，则$F(x)$在$(0,1)$上单调递增，可得$F(x)>F(0)=0$，所以$x>\sin x$，$x\in(0,1)$；构建$G(x)=\sin x-(x-x^{2})=x-x^{2}+\sin x$，$x\in(0,1)$，则$G'(x)=2x-1+\cos x$，$x\in(0,1)$，构建$g(x)=G'(x)$，$x\in(0,1)$，则$g'(x)=2-\sin x>0$对$\forall x\in(0,1)$恒成立，则$g(x)$在$(0,1)$上单调递增，可得$g(x)>g(0)=0$，即$G'(x)>0$对$\forall x\in(0,1)$恒成立，则$G(x)$在$(0,1)$上单调递增，可得$G(x)>G(0)=0$，所以$\sin x>x-x^{2}$，$x\in(0,1)$；综上所述：$x-x^{2}<\sin x<x$.
（2）令$1-x^{2}>0$，解得$-1<x<1$，即函数$f(x)$的定义域为$(-1,1)$，若$a=0$，则$f(x)=-\ln(1-x^{2})$，$x\in(-1,1)$，因为$y=-\ln u$在定义域内单调递减，$y=1-x^{2}$在$(-1,0)$上单调递增，在$(0,1)$上单调递减，则$f(x)=-\ln(1-x^{2})$在$(-1,0)$上单调递减，在$(0,1)$上单调递增，故$x=0$是$f(x)$的极小值点，不合题意，所以$a\neq0$.
当$a\neq0$时，令$b=|a|>0$，因为$f(x)=\cos ax-\ln(1-x^{2})=\cos(|a|x)-\ln(1-x^{2})=\cos bx-\ln(1-x^{2})$，且$f(-x)=\cos(-bx)-\ln[1-(-x)^{2}]=\cos bx-\ln(1-x^{2})=f(x)$，所以函数$f(x)$在定义域内为偶函数，由题意可得：$f'(x)=-b\sin bx-\dfrac{2x}{x^{2}-1}$，$x\in(-1,1)$.
（i）当$0<b^{2}\le 2$时，取$m=\min\{\dfrac{1}{b},1\}$，$x\in(0,m)$，则$bx\in(0,1)$，由（1）可得$f'(x)=-b\sin(bx)-\dfrac{2x}{x^{2}-1}>-b^{2}x-\dfrac{2x}{x^{2}-1}=\dfrac{x(b^{2}x^{2}+2-b^{2})}{1-x^{2}}$，且$b^{2}x^{2}>0$，$2-b^{2}\ge0$，$1-x^{2}>0$，所以$f'(x)>\dfrac{x(b^{2}x^{2}+2-b^{2})}{1-x^{2}}>0$，即当$x\in(0,m)\subset(0,1)$时，$f'(x)>0$，则$f(x)$在$(0,m)$上单调递增，结合偶函数的对称性可知：$f(x)$在$(-m,0)$上单调递减，所以$x=0$是$f(x)$的极小值点，不合题意；
（ii）当$b^{2}>2$时，取$x\in(0,\dfrac{1}{b})\subset(0,1)$，则$bx\in(0,1)$，由（1）可得$f'(x)=-b\sin(bx)-\dfrac{2x}{x^{2}-1}<-b(bx-b^{2}x^{2})-\dfrac{2x}{x^{2}-1}=\dfrac{x(-b^{3}x^{3}+b^{2}x^{2}+b^{3}x+2-b^{2})}{1-x^{2}}$，构建$h(x)=-b^{3}x^{3}+b^{2}x^{2}+b^{3}x+2-b^{2}$，$x\in(0,\dfrac{1}{b})$，则$h'(x)=-3b^{3}x^{2}+2b^{2}x+b^{3}$，$x\in(0,\dfrac{1}{b})$，且$h'(0)=b^{3}>0$，$h'(\dfrac{1}{b})=b^{3}-b>0$，则$h'(x)>0$对$\forall x\in(0,\dfrac{1}{b})$恒成立，可知$h(x)$在$(0,\dfrac{1}{b})$上单调递增，且$h(0)=2-b^{2}<0$，$h(\dfrac{1}{b})=2>0$，所以$h(x)$在$(0,\dfrac{1}{b})$内存在唯一的零点$n\in(0,\dfrac{1}{b})$，当$x\in(0,n)$时，则$h(x)<0$，且$x>0$，$1-x^{2}>0$，则$f'(x)<\dfrac{x}{1-x^{2}}(-b^{3}x^{3}+b^{2}x^{2}+b^{3}x+2-b^{2})<0$，即当$x\in(0,n)\subset(0,1)$时，$f'(x)<0$，则$f(x)$在$(0,n)$上单调递减，结合偶函数的对称性可知：$f(x)$在$(-n,0)$上单调递增，所以$x=0$是$f(x)$的极大值点，符合题意；
综上所述：$b^{2}>2$，即$a^{2}>2$，解得$a>\sqrt{2}$或$a<-\sqrt{2}$，故$a$的取值范围为$(-\infty,-\sqrt{2})\cup(\sqrt{2},+\infty)$.
\end{solution}
\end{question}

